\section{Algorithms}
\label{sec:compute}

Four computational problems follow, treated in the order a user meets them.
Compute the answer on the shortlist as given. Then ask what it would be if the
normalisation bounds were different, first by sampling and then, if the budget
allows, with a guarantee holding over a whole continuum of bounds. Finally, when
the alternatives are not a list but the solutions of a model, compute the same
selection without generating that list.

\subsection{The finite problem}

The exact procedure fixes the active set, computes the bounds and anchors,
removes degenerate probes while warning about any criterion that loses its last
separating probe, evaluates the disappointment matrix, sorts each row worst
first, and returns the lexicographically minimal class. Its cost is dominated by
evaluating the family, not by comparing alternatives. Building the matrix costs
$O\big(N\sum_q T_q\big)$ where $T_q$ is the cost of one probe, and sorting the
rows costs $O(NK\log K)$. For the canonical core, all $K=m+2$ probes can be evaluated for one alternative in $O(m)$ operations combined, and sorting them costs $O(m\log m)$. Thus, the
whole selection is $O(Nm\log m)$. In the reference implementation (Python 3.11, Apple M3 Max, macOS 14.5), this resolves a hundred thousand candidates on
ten criteria in a median of $200$ milliseconds (IQR: 3 ms), single threaded, over a hundred
repetitions. The measured interval covers normalisation, probe construction and
evaluation, anchor computation, sorting and comparison. It excludes generating
and dominance-filtering the candidate set, which happens upstream. The number matters not because selection
is a bottleneck, but because the analyses below re-run this procedure thousands
of times.

\subsection{Uncertain bounds: sampling with a guarantee}

Write $\mathcal B$ for a box of admissible bound vectors and $W_0(b)$ for the
exact winner class under bounds $b$. The two objects of interest are the
alternatives that win under \emph{some} admissible bound and those that win
under \emph{all} of them,
\begin{equation}
S^{\mathrm{pos}}=\bigcup_{b\in\mathcal B}W_0(b),
\qquad
S^{\mathrm{nec}}=\bigcap_{b\in\mathcal B}W_0(b).
\label{eq:interval}
\end{equation}
Every experiment here uses one model for $\mathcal B$, set by a single level
$p\ge0$ and applied multiplicatively to each nominal bound, independently across
criteria. The ideal $\ideal_i$ is multiplied by $1-u'_i$ with
$u'_i\sim\mathrm{Unif}(0,p)$ and the nadir $\nadir_i$ by $1+u_i$ with
$u_i\sim\mathrm{Unif}(-p,p)$, the nadir floored above the ideal so every draw has
positive range. Because perturbed nadir bounds can fall within the observed criterion range, they act as scale parameters rather than strict extremes. If strict bracketing is required, the perturbation can be constrained to the domain bounding the observed candidates. The asymmetry is deliberate: for a cost criterion the true best
attainable value can lie below the observed one, but the observed one is
attained. The certified mode of the next subsection
takes the deterministic hull of this support, the box $\mathcal B(p)$. A ``five
per cent box'' means $p=0.05$ in this sense, not a fraction of the criterion
range. Two scope conditions come with the multiplicative form. It presumes bounds of
constant sign, which holds here because every criterion is a positive cost or
quantity. And the sampling distribution is part of the declared model, not a
neutral default, so every frequency reported under it is conditional on it.

Sampling $M$ draws gives a lower approximation of $S^{\mathrm{pos}}$ with a
one-sided guarantee. This is a pointwise guarantee: for any fixed alternative whose winning region has probability at least $\alpha$, the probability that $M$ independent draws all miss it is at most $(1-\alpha)^M$. By the union bound, to simultaneously observe all such alternatives with confidence at least $1-\delta$, a sample size of $M \ge \log(N/\delta)/(-\log(1-\alpha))$ is sufficient. That turns ``we sampled a thousand draws'' into a statement
with a failure probability attached. It says nothing about
winning regions of small probability, and nothing at all about regions of
probability zero, which is the gap the next subsection closes. The direction of
the guarantee also has to be respected. The sampled union is an \emph{inner}
approximation of $S^{\mathrm{pos}}$, so finding an alternative there proves it
can win. The sampled intersection is an \emph{outer} approximation of
$S^{\mathrm{nec}}$, so it can only be too large, and it must never be read as a
certificate that an alternative wins under every admissible bound.

\subsection{Uncertain bounds: certification without sampling}

The certified mode replaces sampling by interval arithmetic and branch and
bound. Single-criterion disappointments do not depend on the bounds, so only aggregate probes have to be
enclosed at all, and the problem is much smaller than its $2m$ nominal
dimensions suggest.

Fix a sub-box $\beta\subseteq\mathcal B$. For fixed $x$ and $i$, the normalised
value $r_i(x;b)$ is linear-fractional in the two bound variables with a
denominator bounded away from zero, so its extremes over the rectangle $\beta$
occur at vertices and can be evaluated exactly. Each canonical probe is monotone
in $r$, so an enclosure of $q(r(x;b))$ follows coordinatewise, and for the
maximum probe it is not merely sound but exact, because the coordinate blocks
are independent under a Cartesian box. The disappointment is a ratio of three
quantities all driven by the same $b$, and the enclosure ignores that shared
dependence, which makes it sound but not tight. Every operation is rounded
outward, one step of the machine's next-representable-value function per
elementary operation, so the enclosures are valid in floating point and not only
in exact arithmetic. If the enclosed active range of a probe fails to be
positive, the trivial enclosure $[0,1]$ is used and the box is subdivided rather
than resolved.

\begin{lemma}[Componentwise monotonicity of $\Psi$]
If $D(x) \leq D(y)$ coordinatewise, then $\Psi(x) \leq \Psi(y)$ coordinatewise.
\end{lemma}
\noindent\emph{Proof.} By monotonicity of empirical CVaR, $D(x) \leq D(y) \implies T_\alpha(x) \leq T_\alpha(y)$ for all $\alpha$. Since $B_\delta$ is a nondecreasing function, $C_\alpha(x) \leq C_\alpha(y)$. Sorting is also monotone, so $D^\downarrow(x) \leq D^\downarrow(y)$. Thus every coordinate of $\Psi(x)$ is bounded by the corresponding coordinate of $\Psi(y)$. $\square$

\begin{proposition}[Soundness of the certified test]
\label{prop:certsound}
Write $\reg^{L}(\cdot;\beta)$ and $\reg^{U}(\cdot;\beta)$ for the coordinatewise
lower and upper enclosures over $\beta$. If
$\Psi\!\big(\reg^{U}(x;\beta)\big)\lex
 \Psi\!\big(\reg^{L}(y;\beta)\big)$ strictly, then
$x\prec_{\mathrm{ReLexTail}}y$ for every $b\in\beta$.
\end{proposition}

\noindent\emph{Proof.} Over any $b \in \beta$, the true disappointment profile satisfies $\reg(y;b) \geq \reg^{L}(y;\beta)$ and $\reg(x;b) \leq \reg^{U}(x;\beta)$ coordinatewise. By the Lemma, this implies $\Psi(\reg(y;b)) \geq \Psi(\reg^{L}(y;\beta))$ and $\Psi(\reg(x;b)) \leq \Psi(\reg^{U}(x;\beta))$ coordinatewise. The lexicographic extension of a coordinatewise inequality is a weak inequality. Thus $\Psi(x;b) \leq_{\mathrm{lex}} \Psi(\reg^{U}(x;\beta)) <_{\mathrm{lex}} \Psi(\reg^{L}(y;\beta)) \leq_{\mathrm{lex}} \Psi(y;b)$, establishing strict preference for $x$. $\square$

A negative answer from this test is therefore trustworthy in a way no sample can
be: an eliminated candidate provably cannot win anywhere in the box. The search
maintains a stack of sub-boxes, each carrying the set of candidates its parent
could not eliminate. On each box it eliminates every candidate that some rival
beats throughout, and it declares the box \emph{resolved} when one candidate
beats every rival throughout. A resolved box contributes its sole winner to the
inner set and to the outer set. An unresolved box, whether depth-capped or left
on the stack when the budget runs out, contributes its inherited survivor set to
the outer set and its volume to the unresolved total.

Two details decide whether this is correct, and both are easy to get wrong. The
outer set is accumulated over terminal leaves by \emph{union}, never by deleting
from a running global set. A candidate eliminated on one sub-box may win on
another, so if $a$ wins throughout $\beta_1$ and $b$ throughout $\beta_2$ the
possible set is $\{a,b\}$, whereas subtracting each local elimination in turn
would remove first $b$ and then $a$ and return the empty set. And resolution
uses the direct test, that some $x^\star$ beats every rival, rather than the
weaker observation that $x^\star$ has no rival certified to beat it, which would
need a separate acyclicity argument. Because elimination is monotone under
subdivision, a parent's survivor set is sound for any descendant, including one
never processed, which is what makes budget exhaustion safe. The archive
contains a regression test that fails if the union is replaced by subtraction.

The procedure returns $\mathrm{In}\subseteq S^{\mathrm{pos}}\subseteq\mathrm{Out}$
and an unresolved volume fraction $u$. When $u=0$ and every resolved box has the
same sole winner, that winner is additionally certified necessary. This test
detects a necessary winner but not a necessary \emph{class}, so the output is a
certified subset of $S^{\mathrm{nec}}$ rather than the set itself unless
witnesses supply the rest. Two
convergence caveats are permanent rather than budgetary. Exact recovery fails on
persistent tie surfaces, where no sole winner exists to certify. And the inner
set can miss an alternative whose winning region has empty interior, since such
an alternative is never the sole winner on a box of positive volume. Its
membership has to be established by an explicit witness, that is by exhibiting
one admissible bound vector at which it wins.

\paragraph{Two radii, not one}
Bisecting the level $p$ under the same test gives a lower bound on the stability
radius, and the paper keeps two quantities apart. The \emph{true radius} of the
nominal winner is
$\rho^\star=\sup\{p\ge0: W_0(b)=\{x^\star\}\ \forall b\in\mathcal B(p)\}$,
a property of the data. The \emph{certification frontier}
$\underline{\rho}_{B,d}$ is the largest $p$ at which the procedure proves
uniqueness under box budget $B$ and depth cap $d$, and it depends on the
enclosures and the budget as well as on the data. A witnessed reversal, an
explicit $b$ at which some rival wins, gives an upper bound $\overline{\rho}$.
Soundness gives
\begin{equation}
\underline{\rho}_{B,d}\ \le\ \rho^\star\ \le\ \overline{\rho},
\label{eq:radbracket}
\end{equation}
and both inequalities can be strict. Failure to certify above the frontier may
mean loose enclosures or an exhausted budget rather than genuine instability, so
the frontier is never reported as the radius. Raising the budget can only raise
it, and no finite budget is guaranteed to close the bracket.

\paragraph{Does it actually work?}
Soundness is a theorem, but a sound procedure whose enclosures are too loose to
prove anything would still be useless, so we test it under a dense-grid stress test.
On low-dimensional instances the box has only $2m$ coordinates and can be swept
exhaustively: we evaluate the exact rule at every point of a dense product grid
over $\mathcal B(p)$ and collect the winners it produces as a set
$S_{\mathrm{grid}}$. Every grid point is an admissible bound vector, so
$S_{\mathrm{grid}}\subseteq S^{\mathrm{pos}}$ by construction, and the inclusion
$\mathrm{In}\subseteq S_{\mathrm{grid}}\subseteq\mathrm{Out}$ is a real test of
the outer set: an alternative that wins somewhere on the grid but is missing
from $\mathrm{Out}$ would be a soundness violation.

Table~\ref{tab:truth} reports the outcome over eleven instances at three levels.
The inclusion holds in all thirty-three cases with no violation, and monotone
tightening holds everywhere: on every row the outer set and the unresolved
fraction are non-increasing in the budget, as terminal-leaf aggregation
requires. Three rows reach exact recovery, the unresolved volume falling to zero
with the inner and outer sets coinciding as a singleton, at budgets of $250$,
$250$ and $4{,}000$ boxes. The rest show the failure modes rather than hiding
them. A constructed instance with two identical criterion rows ties at every
grid point, so no sole winner exists anywhere and the inner set stays empty:
correct behaviour, not failure. A near-degenerate instance whose second
criterion spans a range of $0.003$ triggers the denominator fallback on $81$ to
$83$ of the $128$ pairs probed, that is $64$ sub-boxes times the two aggregate
probes. And on most rows the grid shows two or three winners while the outer set is
still the whole candidate set, the visible cost of ignoring the shared
dependence in the enclosure.

\begin{table}[t]
\centering\small
\setlength{\tabcolsep}{3.5pt}
\caption{Dense-grid stress test of the certified outer set, eleven instances at three levels.
$S_{\mathrm{grid}}$ is the winner set found by sweeping the bound box
exhaustively. The last column is the smallest budget on the ladder
$250,1{,}000,4{,}000,16{,}000$ at which the inner and outer sets coincide with
zero unresolved volume, and is blank where no rung achieved it. Full per-row
output is in the archive.}
\label{tab:truth}
\begin{tabular}{@{}P{2.0cm}rrP{1.15cm}P{1.35cm}@{}}
\toprule
Behaviour & rows & viol. & grid ties & $B^{\mathrm{ex}}$\\
\midrule
Exact recovery & 3 & 0 & $0\%$ & $250$--$4{,}000$\\
Dependency slack & 24 & 0 & $0\%$ & ---\\
Persistent tie & 3 & 0 & $100\%$ & ---\\
Near-degenerate & 3 & 0 & $0.5$--$0.8\%$ & ---\\
\midrule
\textbf{All} & \textbf{33} & \textbf{0} & & \\
\bottomrule
\end{tabular}
\end{table}

Two limits are worth naming. The grid is dense but finite, $23$ points per
coordinate at $m=2$ and $8$ at $m=3$, so a winning region smaller than the grid
step would be invisible to it. And the certified mode reaches only small boxes
and modest criterion counts, and elsewhere the honest output is the bracket with its
unresolved volume attached.

\paragraph{What certification costs}
Two things govern the practical reach and they are not the same. The box count
is the visible budget, but the cost per box grows with the candidate set,
because certifying a box means showing that every candidate is beaten
throughout. Measured cost per box rises from $0.87$ milliseconds at $N=20$ to
$8.11$ at $N=100$, nearly a factor of ten for a fivefold increase in $N$. Pure quadratic
growth would give twenty-five, so the realised cost sits between linear and
quadratic, because the elimination test stops at the first rival that
dominates. A hundred-alternative shortlist is \emph{processed} in about eight seconds at a
thousand-box budget, which is not the same as certified: that run left $93\%$ of
the box unresolved and returned the whole shortlist as its outer set. Time is
cheap here and resolution is not, which is the point of the previous paragraph.
The validated interval arithmetic costs roughly a fifth of the running time and
changes no certified set, box count or unresolved fraction: switching from plain
floating point to outward-rounded operations reproduces every certified number
here to the digit.

Figure~\ref{fig:certdecay} shows how the unresolved volume falls as the budget
grows, and the shape of that curve is the practical message. Quadrupling the
budget from $250$ to $1{,}000$ sub-boxes takes the unresolved volume from
$0.766$ to $0.569$, and doubling again to $2{,}000$ takes it to $0.097$. At a
nominal budget of $4{,}000$ the procedure terminates inside its budget, using
$2{,}277$ boxes and leaving no unresolved leaf. That last step is where the
guarantee stops being merely sound and becomes sharp: throughout the
under-resolved regime the outer set is the full candidate set, so the procedure
proves nothing about who wins, only that its answer is not wrong, and the moment
the unresolved volume reaches zero the outer set collapses from fifteen
alternatives to one. Soundness holds at every budget, sharpness arrives all at once. The unresolved fraction, not the running time, is the number that tells an
analyst which regime a run is in.

\begin{figure}[t]
\centering
\includegraphics[width=\linewidth]{fig_certified_decay.pdf}
\caption{Spending more effort shrinks the part of the bound box left undecided.
What grows with the budget is not the validity of the answer, which holds
throughout, but how much of the box the procedure manages to settle.}
\label{fig:certdecay}
\end{figure}

\subsection{When the alternatives are not a list}

Sometimes the candidate set is not given but implied by a model, and enumerating
it is the expensive part. If the feasible set $\X$ is a nonempty bounded
polyhedron, the criteria are affine, and the probes are canonical,
the ReLexTail selection can be computed directly as a sequence of mixed-integer linear programmes (MILPs).

The ReLexTail profile $\Psi(x)$ is evaluated sequentially.
For a risk functional $R_\ell(x)$, the integer category variable $z_\ell$ is defined by:
\[ R_\ell(x) \leq \delta_\ell z_\ell, \qquad z_\ell \in \mathbb{Z}_{\geq 0}. \]
Minimising $z_\ell$ enforces $z_\ell = \lceil R_\ell(x)/\delta_\ell \rceil$.
For the maximum $M(x)$, an epigraph variable $v_M \geq D_q(x)$ enforces $M(x) \leq \delta_M z_M$.
For the empirical CVaR $T_\alpha(x)$, the variational identity of \citet{rockafellar2000} is used:
\[ s_{\alpha q} \geq D_q(x) - t_\alpha, \quad s_{\alpha q} \geq 0, \]
\[ t_\alpha + \frac{1}{\alpha K} \sum_q s_{\alpha q} \leq \delta_\alpha z_\alpha. \]

The ReLexTail profile is minimised sequentially through $K+7$ stages:
\begin{enumerate}
    \item \textbf{Category stages (MILP):} Minimise $z_M$; constrain $z_M=z_M^\star$, minimise $z_{25}$; fix previous, minimise $z_{50}$; fix previous, minimise $z_{100}$. If $\X$ is continuous, these stages are MILPs because $z_\ell \in \mathbb{Z}$.
    \item \textbf{Exact tail stages (LP/MILP):} Minimise $T_{25}$; minimise $T_{50}$; minimise $T_{100}$; minimise $M$.
    \item \textbf{Terminal refinement (LP/MILP):} Since $M=D_{[1]}$ is already minimised, we preserve all previous optima and minimise the sum of the two largest disappointments, continuing for $p=3,\ldots,K$. This requires $K-1$ exact-leximax stages.
\end{enumerate}

\begin{proposition}[Solver call count]
\label{prop:count}
The canonical implementation requires at most $4m+12$ optimisation stages under the stated reuse convention.
\end{proposition}

\noindent\emph{Proof.} Finding criterion bounds requires $2m$ programmes. The grand-mean and grand-maximum anchor preparations require $2$ and $m+1$ programmes, respectively. The $K+7$ ReLexTail stages require $(m+2)+7 = m+9$ stages. The total is $2m + 2 + (m+1) + (m+9) = 4m+12$. $\square$

\begin{theorem}[Direct ReLexTail computation]
\label{thm:continuous}
Under bounded polyhedral $\X$, affine criteria, canonical probes, and positive resolutions, the optimal ReLexTail profile $\min_{\mathrm{lex}}\Psi(x)$ can be computed by a finite sequence of mixed-integer and continuous linear programmes, characterising the optimal solution set through the final stage-optimality constraints.
\end{theorem}

\noindent\emph{Proof.} By formulation, the stage variables correctly represent the category and tail components of $\Psi(x)$. The sequential imposition of optimal values exactly mirrors the lexicographic priority of $\Psi$. Boundedness ensures finiteness and attainment. $\square$

Three qualifications belong with this result and are usually left out. First,
the optimal set may be infinite, and it need not be a face of $\X$: minimising a
convex piecewise-linear function over a polytope can select a relative interior
point, as $\min\max\{x,1-x\}$ over $[0,1]$ does at $x=1/2$. The procedure
therefore characterises a polyhedral optimal subset, and a declared secondary
rule, itself recorded, picks a representative when one point is required.

Second, the word ``exactly'' holds under exact optimisation. A floating-point solver
returns optima up to its tolerance, and imposing a previously reported optimum
as a hard constraint on a later stage is not safe: the stage can become
infeasible, or can exclude points that were optimal in exact arithmetic. The
implementation therefore imposes the carried optimum inflated by
$\epsilon_{p'}=\texttt{tol}\cdot(1+|T^\star_{p'}|)$ with
$\texttt{tol}=10^{-6}$, recorded in the audit record with the solver identity.
The inflation is relative by design, because stage optima grow with the stage,
being sums of order statistics, so a fixed absolute tolerance that is adequate
at the first stage is too tight at the last and shows up as a spurious
infeasibility. The price is that the recovered set is optimal for an
$\epsilon$-relaxed sequence and can be strictly larger than the exact one. This
is a solver-tolerance matter, unrelated to the interval certification above.

Third, the anchors here are computed over $\X$, whereas the finite procedure
computes them over the filtered candidate set. Those anchors differ, so the two
versions need not agree. Avoiding enumeration of the efficient set is not the
same as reproducing the answer one would get from it, and which convention is
intended has to be declared.

\paragraph{What the computational study does and does not show}
The scheme is implemented in full and measured against a comparator that draws
$400$ weight vectors, solves a weighted-sum programme for each, and applies the
leximax comparison to the resulting points. Both rules use Gurobi 11.0.2. Figure~\ref{fig:continuous} reports
the cost over fifty instances. The realised solver-call count is $6m+5$ at every size, and it matches the
scheme term by term rather than merely respecting its order. The direct
formulation uses about an order of magnitude fewer solver calls, $53$ against
$443$ at $m=8$, and runs five times faster. The comparator budget is held fixed at
four hundred weight vectors at every size, so the comparison does not test how
that budget would have to grow with $m$ to keep the sampled set representative.
The claim is that the direct cost is linear in $m$ against a fixed comparator
budget, not that the gap provably widens.

We do not present the accompanying quality comparison as validation of
optimality, and it would be misleading to do so. In all fifty instances the direct solution was lexicographically no worse than
the best sampled point, but that is close to tautological. The comparator searches
a finite subset of the same feasible set, and weighted sums cannot reach
unsupported efficient points at all, so the direct solution must be no
worse. Calling four hundred weighted sums
an enumeration of the efficient set would be wrong. Two gaps are open rather
than hidden: an independent check against exact vertex enumeration on small
problems would test optimality properly and is not attempted, and
Theorem~\ref{thm:continuous} covers mixed-integer sets while the study exercises
polytopes only, so that branch is proved and not benchmarked.

\begin{figure}[t]
\centering
\includegraphics[width=\linewidth]{fig_continuous_cost.pdf}
\caption{Solving the selection directly over the polytope costs a small and
predictable number of linear programmes, exactly $6m+5$ for the canonical
family, while weighted-sum sampling costs hundreds whatever the criterion count.
The vertical scale is logarithmic.}
\label{fig:continuous}
\end{figure}

\section{Experiments}
\label{sec:experiments}

The empirical study answers whether the categorical resolution protection of ReLexTail succeeds at defending against macroscopic worst-case deterioration while improving stability over exact leximax, and isolates which component of the rule drives the observed behaviour.

We evaluate $500$ candidate sets from six generators: knapsack and job-shop instances from combinatorial and $\varepsilon$-constraint procedures, WFG2 \citep{huband2006} and DTLZ \citep{deb2002} sampled by Latin hypercube \citep{mckay1979}, and Energy \citep{uci_energy} and Concrete \citep{uci_concrete} surrogate matrices. Sets reach $N=500$ points and $m=15$ criteria, and none is an exactly enumerated Pareto front, which is why we call them candidate sets throughout. The declared family is the canonical core \eqref{eq:qcan}. Table~\ref{tab:protocol} records the protocol.

\begin{table}[t]
\centering\small
\setlength{\tabcolsep}{3.5pt}
\caption{Protocol of the sampled and certified experiments.}
\label{tab:protocol}
\begin{tabular}{@{}P{2.1cm}rrP{1.0cm}P{1.45cm}@{}}
\toprule
Experiment & Inst. & Sds & Draws & Levels\\
\midrule
Selection, divergence & $500$ & $1$ & --- & ---\\
Point fragility & $500$ & $1$ & $60$/lvl & $5,10,20\%$\\
Decomposition & $5{,}000$ & $10$ & $60$/lvl & $5,10,20\%$\\
Supplier flip curve & $1$ & $1$ & $300$/lvl & $0$--$50\%$\\
Certified supp. & $1$ & --- & $40{,}000$ boxes & $2$--$20\%$\\
Known-truth & $11$ & $1$ & none & $2,5,10\%$\\
Polyhedral study & $50$ & $1$ & --- & ---\\
\bottomrule
\end{tabular}
\end{table}

\subsection{RQ1: Representation effect}
Does active-range re-anchoring reduce sensitivity relative to raw normalization?
A naive comparison of a lexicographic rule against a raw-normalised achievement scalarisation confounds the richness of the family, the re-anchoring, and the order itself. Re-anchoring carries the entire initial stability gap. Moving from raw normalisation to re-anchored disappointments takes the class flip rate from $75.3\%$ down to about $44.5\%$ at a $20\%$ bound perturbation level. This holds regardless of whether the profile is ordered lexicographically or collapsed by a scalarisation. 

\subsection{RQ2: Resolution-order effect}
Does ReLexTail improve loss, stability, or class size relative to exact LexPR?
Exact LexPR is highly fragile because it gives infinite priority to microscopic differences in the absolute worst probe. Figure~\ref{fig:quality_stability} displays the Quality-Stability frontier. By mapping upper-tail empirical CVaR into resolution categories, ReLexTail strictly dominates exact LexPR on the frontier: it halves the hidden-preference tail loss while simultaneously reducing the class flip rate.

\begin{figure}[t]
\centering
\includegraphics[width=\linewidth]{quality_stability_frontier.pdf}
\caption{The Quality-Stability frontier. Points show the tail hidden-preference loss against the point flip rate for the comparator methods and ReLexTail with varying resolution thresholds ($\delta$).}
\label{fig:quality_stability}
\end{figure}

\subsection{RQ3: Resolution sensitivity}
How does the resolution threshold $\delta$ move the frontier?
Table~\ref{tab:decomp} shows that as $\delta$ increases (coarser resolution), stability improves but at the cost of higher tail loss, because the rule becomes more willing to accept moderate worst-case deterioration in exchange for better exact-tail refinements. A threshold of $\delta=0.01$ provides the strongest compromise, significantly outperforming strict scalar methods.

\begin{table}[t]
\centering\small
\caption{Decomposition of the fragility gap over the pooled population, $60$ bound draws per level. ReLexTail variants use different resolution thresholds.}
\label{tab:decomp}
% Auto-generated by scripts/run_symmetrised.py. Do not edit.
\footnotesize
\begin{tabular}{@{}lrrrrr@{}}
\toprule
Rule & \multicolumn{3}{c}{Flip rate} & Tail & Agrees\\
\cmidrule(lr){2-4}
 & $5\%$ & $10\%$ & $20\%$ & loss & w/ LexPR\\
\midrule
ASF$_{\mathrm{raw}}$ & 36.4\% & 55.5\% & 75.1\% & 0.584 & ---\\
\textbf{LexPR} & 18.3\% & 29.9\% & 44.8\% & 0.597 & ---\\
ASF$_K$ & 17.9\% & 29.5\% & 44.5\% & 0.597 & 98.0\%\\
MMR$_K$ & 18.4\% & 30.0\% & 44.8\% & 0.597 & 97.6\%\\
LexPR$_1$ & 0.0\% & 0.0\% & 0.0\% & 0.584 & ---\\
ASF$_1$ & 0.0\% & 0.0\% & 0.0\% & 0.584 & ---\\
\bottomrule
\end{tabular}

\end{table}

\subsection{RQ4: Structural attribution}
Which components matter? The stability comes from normalising each declared question on the range that question actually spans (RQ1). The quality and explanation power come from evaluating the empirical CVaR across canonical probes (RQ2). The category discretisation is what prevents arbitrarily small differences from causing catastrophic point fragility.

\subsection{RQ5: Generator heterogeneity}
Are effects consistent across families? Divergence is near total on knapsack instances and near zero on WFG2. Figure~\ref{fig:divergence} illustrates this. Any statement of the form ``ReLexTail is less fragile'' must be qualified by generator family, although the relative ranking of rules is consistent.

\begin{figure}[t]
\centering
\includegraphics[width=\linewidth]{fig_divergence.pdf}
\caption{How often the two rules pick different alternatives, by generator, on the pooled population.}
\label{fig:divergence}
\end{figure}

\subsection{RQ6: Certification}
Can possible-winner sets be resolved under the discontinuous category map?
The interval mode replaces sampling by interval arithmetic. Because the category function $B_\delta(z) = \lceil z/\delta \rceil$ is monotonically non-decreasing, the coordinatewise enclosures of empirical CVaR propagate soundly through the ReLexTail profile. Figure~\ref{fig:certdecay} shows the unresolved volume falling as the box budget grows, confirming that exact certification remains practical despite the discontinuous formulation.

\begin{figure}[t]
\centering
\includegraphics[width=\linewidth]{fig_certified_decay.pdf}
\caption{Spending more effort shrinks the part of the bound box left undecided.}
\label{fig:certdecay}
\end{figure}

\subsection{RQ7: Direct optimisation}
Does the MILP reproduce enumerated ReLexTail optima?
The mixed-integer linear programme introduced in Section~\ref{sec:compute} correctly recovers the optimal ReLexTail profile. As established in Proposition~\ref{prop:count}, the number of solver stages scales as $4m+12$, a small predictable sequence. Figure~\ref{fig:continuous} shows the cost over fifty polyhedral instances, confirming the theoretical bound.

\begin{figure}[t]
\centering
\includegraphics[width=\linewidth]{fig_continuous_cost.pdf}
\caption{Solving the selection directly over the polytope costs exactly $4m+12$ stages for the canonical family.}
\label{fig:continuous}
\end{figure}

\subsection{One instance, end to end}
We demonstrate the exact diagnostic outputs on a ten-supplier instance evaluated on seven cost criteria. The ReLexTail decision profile for the recommended supplier (S7) is shown in Figure~\ref{fig:relex_profile}. 

\begin{figure}[t]
\centering
\includegraphics[width=\linewidth]{relex_decision_profile.pdf}
\caption{The ReLexTail decision profile. Panel A shows the labelled probe disappointments. Panel B shows the tail spectrum $T_\alpha(x)$ for varying $\alpha$. Panel C shows the resolution categories.}
\label{fig:relex_profile}
\end{figure}

By inspecting the category profiles, we find the decision is decisive at the first category tier (the maximum), with exact decisive margin $\Delta_{\text{exact}} = \supReLexDeltaExact$, confirming stability against micro-perturbations. Under interval bounds, the point recommendation is predictably fragile but the interval mode successfully brackets the necessary winners.

\begin{figure}[t]
\centering
\includegraphics[width=\linewidth]{fig_supplier_interval.pdf}
\caption{The point recommendation is fragile and the interval mode says so.}
\label{fig:interval}
\end{figure}



\section{Discussion}
\label{sec:discussion}

ReLexTail is useful when the decision maker cannot commit to precise trade-off
weights but can declare and defend a finite family of monotone scalar questions,
and when the recommendation has to be explained and re-checked afterwards. It is
not the right tool when the objective is average-case utility under a known
preference distribution, when no meaningful monotone probes can be declared or
the family is itself the political dispute, or when the bounds are so uncertain
that the stability class stays large. In that last case the method still gives
the right answer, which is that the data do not support naming one alternative.

Three limitations are structural rather than fixable by more computation. The
rule depends on the candidate set, because active-set normalisation ties the
scale to the observed extremes, which dictates exactly which
edits are safe. Robustness holds relative to the declared family rather than
over all conceivable preferences, as the probe ablation shows. And the construction protects against the worst
case rather than maximising average utility, which the held-out results reflect.

Several threats to validity deserve naming. The supplier case is reproducible
but constructed, so validation on real decision problems is the most important
next step. The held-out preference pool is author-chosen, weights ten families
equally, and contains two that share a power-mean form, so that geometry counts
twice. The certified mode reaches only small boxes and modest criterion counts.
And the polyhedral formulation is proved for canonical probes over polyhedra
with affine criteria, with the mixed-integer branch untested and general convex
probes out of scope.

This paper introduced ReLexTail, a deterministic resolution-aware lexicographic order for auditable multicriteria selection. The takeaway is that the last step of a multicriteria study, turning an efficient set into a defensible choice, can be treated as a rigorous computational problem with formal guarantees. By mapping upper-tail empirical CVaR into discrete categories before breaking ties exactly, ReLexTail protects against macroscopic worst-case deterioration while remaining fully complete and transitive. We showed that this order supports certified possible-winner analysis over a continuum of normalisations without sampling, can be formulated as a sequence of mixed-integer linear programmes over polyhedra to skip front enumeration, and exports a fully replayable audit record. The empirical work confirms that the categorical representation preserves the core stability benefits of active-range normalisation while improving the hidden-preference tail loss compared to exact leximax, effectively upgrading the quality--stability compromise.

\section*{Declarations}

\noindent\textbf{Competing interests.} The author declares no competing
interests.

\noindent\textbf{Funding.} No external funding supported this work.

\noindent\textbf{Ethics approval.} Not applicable. The study uses no human
participants, no animal subjects and no personal data.

\noindent\textbf{Data and code availability.} All data, code and seeds are in
the reproduction archive cited below. The supplier instance is constructed and
is distributed with the archive. No third-party or proprietary data is used.

\noindent\textbf{Relationship to a previous preprint.} An earlier and longer
version of this work was deposited as a preprint (Zenodo record 21037960). The
present paper supersedes it. Two classes of claim made there are withdrawn
rather than restated: an ``auditability trilemma'' framing, and the stronger
axiomatic reading of the profile-order result, which is a representation lemma
relative to a declared family and not a derivation of the leximax order from
primitive postulates. The empirical claims have also been narrowed, as
Table~\ref{tab:claims} records.

\noindent\textbf{Use of large language models.} A large language model was used
as a writing aid: to shorten prose, to check the reported numbers against the
archived outputs, and to draft regression tests. It was not used to generate
results or to design experiments. Every quantity here is computed by the
archived scripts, and the author is responsible for the whole content.

\section*{Reproducibility}

All code, data, seeds and scripts are in the archive \citep{bezoui2026archive},
version~\ArchiveVersion, at
\url{https://doi.org/10.5281/zenodo.21739625} and tagged \ArchiveRelease{} at
\url{https://github.com/MadBezoui/ReLexTail}. One command,
\texttt{python scripts/reproduce\_all.py}, regenerates every number and figure
here from the fixed master seeds. The archive also holds the full proofs, the
algorithm listings with their enclosure formulas, the per-row output of
Table~\ref{tab:truth}, the two supplier witness vectors and the audit
records.
